
\documentclass{article}
\usepackage{colortbl}
\usepackage{makecell}
\usepackage{multirow}
\usepackage{supertabular}

\begin{document}

\newcounter{utterance}

\centering \large Interaction Transcript for game `wordle', experiment `high\_frequency\_words\_no\_clue\_no\_critic', episode 3 with NEW\_Reflect\_Agent.
\vspace{24pt}

{ \footnotesize  \setcounter{utterance}{1}
\setlength{\tabcolsep}{0pt}
\begin{supertabular}{c@{$\;$}|p{.15\linewidth}@{}p{.15\linewidth}p{.15\linewidth}p{.15\linewidth}p{.15\linewidth}p{.15\linewidth}}
    \# & \multicolumn{2}{c}{Player} && \multicolumn{2}{c}{Game Master} \\
    \hline

    \theutterance \stepcounter{utterance}  
    & & & \multicolumn{4}{p{0.6\linewidth}}{
        \cellcolor[rgb]{0.9,0.9,0.9}{
            \makecell[{{p{\linewidth}}}]{
                \texttt{\tiny{[P1$\langle$GM]}}
                \texttt{You are a language wizard who likes to guess words by using the given rules.} \\
\\ 
\texttt{Welcome to Wordle! You have six attempts to guess the target word, a valid English word of five lowercase letters (a{-}z). Please use the tags "explanation:" and "guess:" to provide a concise explanation for each guess.} \\
\\ 
\texttt{For instance, if your guess is "apple", your response should be} \\
\texttt{explanation: this is a common five{-}letter English word, and I am starting my guess with this word.} \\
\texttt{guess: apple} \\
\\ 
\texttt{After each guess, your answer will be validated, and you will receive feedback indicating which letters are correct (green), which letters are correct but in the wrong position (yellow), and which letters are incorrect (red). This feedback can be useful in determining which letters to include or exclude in your next guess.} \\
\\ 
\texttt{For example, the feedback for "apple" might be:} \\
\texttt{guess\_feedback: a<yellow> p<yellow> p<green> l<yellow> e<red>} \\
\\ 
\texttt{The explanation should contain details about how the guess\_feedback is used to arrive at a new guess.} \\
\\ 
\texttt{Do not generate any other text.} \\
\\ 
\texttt{Let's begin with your first guess.} \\
            }
        }
    }
     \\ \\

    \theutterance \stepcounter{utterance}  
    & \multicolumn{4}{p{0.45\linewidth}}{
        \cellcolor[rgb]{0.95,0.95,0.95}{
            \makecell[{{p{\linewidth}}}]{
                \texttt{\tiny{[P1$|internal]}}
                \texttt{Reflections from previous episodes:} \\
\texttt{{-}} \\
\texttt{ Failed. I guessed "crane" after initial feedback, but "crane" had mismatches (c, n red) while a, r, e were green. This suggests the target word likely had "r a e" in positions 2,3,5, with other letters to determine. My guess "crane" was inefficient; a better second guess would have tested more possible letters for positions 1 and 4, using known greens. Next time, prioritize guesses that incorporate confirmed green letters while testing new letters in unknown slots.} \\
\texttt{{-}} \\
\texttt{ Failed. I repeated the same guess multiple times despite identical feedback, wasting attempts. I should have used the feedback to eliminate letters and try new ones, especially focusing on the green 'e' and 'n' in positions 4 and 5. Next time, I will systematically test remaining possible letters for the first three positions instead of stalling.} \\
            }
        }
    }
    & & \\ \\

    \theutterance \stepcounter{utterance}  
    & \multicolumn{4}{p{0.45\linewidth}}{
        \cellcolor[rgb]{0.95,0.95,0.95}{
            \makecell[{{p{\linewidth}}}]{
                \texttt{\tiny{[P1$|internal]}}
                \texttt{{[}attempt 1 – valid{]}} \\
\texttt{explanation: I start with "crane" as my first guess to test common letters and gather initial feedback. This helps identify which letters are in the target word and their potential positions, especially focusing on vowels and frequently used consonants.} \\
\texttt{guess: crane} \\
            }
        }
    }
    & & \\ \\

    \theutterance \stepcounter{utterance}  
    & \multicolumn{4}{p{0.6\linewidth}}{
        \cellcolor[rgb]{0.561,0.737,0.561}{
            \makecell[{{p{\linewidth}}}]{
                \texttt{\tiny{[P1$\rangle$GM]}}
                \texttt{explanation: I start with "crane" as my first guess to test common letters and gather initial feedback. This helps identify which letters are in the target word and their potential positions, especially focusing on vowels and frequently used consonants.} \\
\texttt{guess: crane} \\
            }
        }
    }
    & & \\ \\

    \theutterance \stepcounter{utterance}  
    & & & \multicolumn{2}{p{0.3\linewidth}}{
        \cellcolor[rgb]{0.9,0.9,0.9}{
            \makecell[{{p{\linewidth}}}]{
                \texttt{\tiny{[GM$|$GM]}}
                \texttt{attempts = 1} \\
\texttt{target\_word = click} \\
\texttt{guess = crane} \\
\texttt{guess\_feedback = c<green> r<red> a<red> n<red> e<red>} \\
            }
        }
    }
    & & \\ \\

    \theutterance \stepcounter{utterance}  
    & & & \multicolumn{4}{p{0.6\linewidth}}{
        \cellcolor[rgb]{0.9,0.9,0.9}{
            \makecell[{{p{\linewidth}}}]{
                \texttt{\tiny{[P1$\langle$GM]}}
                \texttt{guess\_feedback: c<green> r<red> a<red> n<red> e<red>} \\
\\ 
\texttt{Provide your response only in this format.} \\
\texttt{explanation: details} \\
\texttt{guess: word} \\
            }
        }
    }
     \\ \\

    \theutterance \stepcounter{utterance}  
    & \multicolumn{4}{p{0.45\linewidth}}{
        \cellcolor[rgb]{0.95,0.95,0.95}{
            \makecell[{{p{\linewidth}}}]{
                \texttt{\tiny{[P1$|internal]}}
                \texttt{{[}attempt 1 – valid{]}} \\
\texttt{explanation: The feedback shows that 'c' is in the correct position (green), while 'r', 'a', 'n', and 'e' are incorrect (red). This means the target word starts with 'c' and does not contain 'r', 'a', 'n', or 'e'. I will use this information to form a new guess that starts with 'c' and avoids the excluded letters.} \\
\texttt{guess: cleft} \\
            }
        }
    }
    & & \\ \\

    \theutterance \stepcounter{utterance}  
    & \multicolumn{4}{p{0.6\linewidth}}{
        \cellcolor[rgb]{0.561,0.737,0.561}{
            \makecell[{{p{\linewidth}}}]{
                \texttt{\tiny{[P1$\rangle$GM]}}
                \texttt{explanation: The feedback shows that 'c' is in the correct position (green), while 'r', 'a', 'n', and 'e' are incorrect (red). This means the target word starts with 'c' and does not contain 'r', 'a', 'n', or 'e'. I will use this information to form a new guess that starts with 'c' and avoids the excluded letters.} \\
\texttt{guess: cleft} \\
            }
        }
    }
    & & \\ \\

    \theutterance \stepcounter{utterance}  
    & & & \multicolumn{2}{p{0.3\linewidth}}{
        \cellcolor[rgb]{0.9,0.9,0.9}{
            \makecell[{{p{\linewidth}}}]{
                \texttt{\tiny{[GM$|$GM]}}
                \texttt{attempts = 2} \\
\texttt{target\_word = click} \\
\texttt{guess = cleft} \\
\texttt{guess\_feedback = c<green> l<green> e<red> f<red> t<red>} \\
            }
        }
    }
    & & \\ \\

    \theutterance \stepcounter{utterance}  
    & & & \multicolumn{4}{p{0.6\linewidth}}{
        \cellcolor[rgb]{0.9,0.9,0.9}{
            \makecell[{{p{\linewidth}}}]{
                \texttt{\tiny{[P1$\langle$GM]}}
                \texttt{guess\_feedback: c<green> l<green> e<red> f<red> t<red>} \\
\\ 
\texttt{Provide your response only in this format.} \\
\texttt{explanation: details} \\
\texttt{guess: word} \\
            }
        }
    }
     \\ \\

    \theutterance \stepcounter{utterance}  
    & \multicolumn{4}{p{0.45\linewidth}}{
        \cellcolor[rgb]{0.95,0.95,0.95}{
            \makecell[{{p{\linewidth}}}]{
                \texttt{\tiny{[P1$|internal]}}
                \texttt{{[}attempt 1 – valid{]}} \\
\texttt{explanation: The first guess revealed that 'c' is in the correct position (green), while 'r', 'a', 'n', and 'e' are incorrect (red). The second guess confirmed that 'c' and 'l' are in the correct positions (green), while 'e', 'f', and 't' are incorrect (red). This means the word starts with 'c', has 'l' in the second position, and does not contain 'r', 'a', 'n', 'e', 'f', or 't'. I will now use this information to form a new guess.} \\
\texttt{guess: clink} \\
            }
        }
    }
    & & \\ \\

    \theutterance \stepcounter{utterance}  
    & \multicolumn{4}{p{0.6\linewidth}}{
        \cellcolor[rgb]{0.561,0.737,0.561}{
            \makecell[{{p{\linewidth}}}]{
                \texttt{\tiny{[P1$\rangle$GM]}}
                \texttt{explanation: The first guess revealed that 'c' is in the correct position (green), while 'r', 'a', 'n', and 'e' are incorrect (red). The second guess confirmed that 'c' and 'l' are in the correct positions (green), while 'e', 'f', and 't' are incorrect (red). This means the word starts with 'c', has 'l' in the second position, and does not contain 'r', 'a', 'n', 'e', 'f', or 't'. I will now use this information to form a new guess.} \\
\texttt{guess: clink} \\
            }
        }
    }
    & & \\ \\

    \theutterance \stepcounter{utterance}  
    & & & \multicolumn{2}{p{0.3\linewidth}}{
        \cellcolor[rgb]{0.9,0.9,0.9}{
            \makecell[{{p{\linewidth}}}]{
                \texttt{\tiny{[GM$|$GM]}}
                \texttt{attempts = 3} \\
\texttt{target\_word = click} \\
\texttt{guess = clink} \\
\texttt{guess\_feedback = c<green> l<green> i<green> n<red> k<green>} \\
            }
        }
    }
    & & \\ \\

    \theutterance \stepcounter{utterance}  
    & & & \multicolumn{4}{p{0.6\linewidth}}{
        \cellcolor[rgb]{0.9,0.9,0.9}{
            \makecell[{{p{\linewidth}}}]{
                \texttt{\tiny{[P1$\langle$GM]}}
                \texttt{guess\_feedback: c<green> l<green> i<green> n<red> k<green>} \\
\\ 
\texttt{Provide your response only in this format.} \\
\texttt{explanation: details} \\
\texttt{guess: word} \\
            }
        }
    }
     \\ \\

    \theutterance \stepcounter{utterance}  
    & \multicolumn{4}{p{0.45\linewidth}}{
        \cellcolor[rgb]{0.95,0.95,0.95}{
            \makecell[{{p{\linewidth}}}]{
                \texttt{\tiny{[P1$|internal]}}
                \texttt{{[}attempt 1 – valid{]}} \\
\texttt{explanation: The first guess revealed that 'c' is in the correct position (first letter). The second guess confirmed 'c' and 'l' are in the correct positions (first and second letters), while 'e', 'f', and 't' are incorrect. The third guess confirmed 'c', 'l', 'i', and 'k' are in the correct positions (first, second, third, and fifth letters), while 'n' is incorrect. The target word must be five letters long, starting with 'c', followed by 'l', then 'i', and ending with 'k'. The fourth letter must be a letter that is not 'n' and fits the remaining constraints.} \\
\texttt{guess: click} \\
            }
        }
    }
    & & \\ \\

    \theutterance \stepcounter{utterance}  
    & \multicolumn{4}{p{0.6\linewidth}}{
        \cellcolor[rgb]{0.561,0.737,0.561}{
            \makecell[{{p{\linewidth}}}]{
                \texttt{\tiny{[P1$\rangle$GM]}}
                \texttt{explanation: The first guess revealed that 'c' is in the correct position (first letter). The second guess confirmed 'c' and 'l' are in the correct positions (first and second letters), while 'e', 'f', and 't' are incorrect. The third guess confirmed 'c', 'l', 'i', and 'k' are in the correct positions (first, second, third, and fifth letters), while 'n' is incorrect. The target word must be five letters long, starting with 'c', followed by 'l', then 'i', and ending with 'k'. The fourth letter must be a letter that is not 'n' and fits the remaining constraints.} \\
\texttt{guess: click} \\
            }
        }
    }
    & & \\ \\

    \theutterance \stepcounter{utterance}  
    & & & \multicolumn{2}{p{0.3\linewidth}}{
        \cellcolor[rgb]{0.9,0.9,0.9}{
            \makecell[{{p{\linewidth}}}]{
                \texttt{\tiny{[GM$|$GM]}}
                \texttt{attempts = 4} \\
\texttt{target\_word = click} \\
\texttt{guess = click} \\
\texttt{guess\_feedback = c<green> l<green> i<green> c<green> k<green>} \\
            }
        }
    }
    & & \\ \\

    \theutterance \stepcounter{utterance}  
    & & & \multicolumn{2}{p{0.3\linewidth}}{
        \cellcolor[rgb]{0.9,0.9,0.9}{
            \makecell[{{p{\linewidth}}}]{
                \texttt{\tiny{[GM$|$GM]}}
                \texttt{game\_result = WIN} \\
            }
        }
    }
    & & \\ \\

\end{supertabular}
}

\end{document}
